 \documentclass{article}
\usepackage{amsmath}
\usepackage{graphicx}
\graphicspath{ {images/} }
\usepackage[spanish]{babel}
\usepackage{bbm}
\usepackage{amsthm}
\usepackage{authblk}
\usepackage{hyperref}
\usepackage[utf8]{inputenc}
\usepackage{mathrsfs}
\usepackage{amsfonts}
\usepackage{amssymb}
\usepackage{enumerate}
\usepackage[left=3cm,right=2cm,top=2cm,bottom=2cm]{geometry}
\usepackage{epsfig}
 \renewcommand{\baselinestretch}{0.93}
 \thispagestyle{empty}
 \pagestyle{empty}
\setlength{\textheight}{53pc} \setlength{\textwidth} {36pc}
\setlength{\topmargin}{-4mm}
\usepackage{multicol}
\newcommand*{\email}[1]{%
    \normalsize\href{mailto:#1}{#1}\par
    }
\setlength{\columnsep}{1cm}
\newcommand{\N}{\mathbb{N}}
\newcommand{\calM}{\cal{M}}
\newcommand{\calP}{\cal{P}}
\newcommand{\F}{\mathbb{F}}
\newcommand{\R}{\mathbb{R}}
\newcommand{\Z}{\mathbb{Z}}
\newcommand{\Q}{\mathbb{Q}}
\newcommand{\C}{\mathbb{C}}
\newcommand{\bbH}{\mathbb{H}}


\theoremstyle{definition}
\newtheorem{ejer}{Ejercicio}
\author{Marcos Morales}
\affil{Universidad Finis Terrae}
\title{Primera Semana\\}



	


\begin{document}


\begin{picture}(400, 75)
   \put(-15,15){\includegraphics[width=5cm]{UFT.png}}
\end{picture}



\begin{center}
\textbf{\Large{Taller 1}} \\
\end{center}

\setlength{\parindent}{0cm}



\begin{center}
\textbf{Cálculo Vectorial}
\end{center}

\hspace{100cm} \\ 


\textbf{Ejercicio 1:} Considere los rectas cuyas ecuaciones vectoriales son

\begin{align*}
 L_1 &: (1,2,0) + t(2,-1,3), \hspace{0.2cm}\text{$ t \in \mathbb{R}$}\\
 L_2 &:  (3,1,3) +\lambda (-1,2,-1), \hspace{0.2cm} \text{$ \lambda \in \mathbb{R}$} \\
 L_3 &:  (5,0,6)+ s(2,-4,2), \hspace{0.2cm}  \text{$ s \in \mathbb{R}$}\\
\end{align*}

a)  Encuentre las ecuaciones generales de cada recta.\\ 

b) Demuestre que las rectas $L_1$ y $L_2$ se intersectan y encuentre el punto de intersección. \\ 

c) Demuestre que las rectas $L_1$ y $L_3$ se intersectan y encuentre el punto de intersección.\\ 

d) Demuestre que las rectas $L_2$ y $L_3$  no se intersectan.\\ 

e) Determine el ángulo entre las rectas $L_1$ y $L_2$ y también el ángulo entre las rectas  $L_1$ y $L_3$. \\ \\ 


\textbf{Ejercicio 2:} Considere los planos dados por las ecuaciones

\begin{align*}
    \Pi_1 &: -3x+5y+2z-3k = 0 \\
    \Pi_2 &: kx+ y-7z +4 =0
\end{align*}

Donde $k$ es una constante. \\

a)  Determine el valor de la constante $k$ para que los planos sean perpendiculares. \\ 

b) Determine la intersección entre los dos planos. \\ \\ 


\textbf{Ejercicio 3:} Considere el campo escalar

\begin{align*}
    f(x,y) = \ln(x^2+y^2-4)
\end{align*}

a) Determine el dominio de $f$. \\

b) Halle las curvas de nivel  $C_0(f)$, $C_1(f)$ y $C_2(f)$. \\

c) Dibuje el dominio de $f$ y las curvas de nivel de la parte $(b)$.








\end{document} 